
\lesson{19/11}
\section{Scattering amplitudes}
The powerful part of QFTs is their predictions of scattering amplitudes
\[
    A=\bra{\text{out}}\ket{\text{in}}
\]
where ``out'' and ``in'' are the final and initial state of a system.

These quantities are obviously interesting in \textit{interacting theories}:
the interacting picture is more complicated than the free picture (as seen in
the IQFT course).

\subsection{Free theory}
The scalar fields of free theory can be expanded as
\[
    \phi(x)=\int_{}^{} \frac{\dd^3p}{(2\pi)^3\sqrt{2E_p}}e^{i\mathbf{p\cdot x}}
    \left(a(t,p)+a^\dagger(t,-p)\right):
\]
consider $t=0$
\[
    \ket{\text{f}}=\int_{}^{} \dd^3x f(0,\mathbf{x})\phi(0,\mathbf{x})\ket{0}=
    \int_{}^{} \dd^3x f(0,\mathbf{x})\int\frac{\dd^3p}{(2\pi)^3\sqrt{2E_p}}e^{i
    \mathbf{p\cdot x}}a^\dagger(0,-p)\ket{0}:=\int_{}^{} \frac{\dd^3p}{2E_p}
    \tilde f(\mathbf{p})\ket{\mathbf{p}}
\]
where
\begin{gather*}
    \ket{\mathbf{p}}:=\sqrt{2E_p}a^\dagger(0,\mathbf{-p})\ket{0},\\
    \tilde f(\mathbf{p})=\int_{}^{} \frac{\dd^3x}{(2\pi)^{\frac{3}{2}}}f(
    \mathbf{x})e^{i\mathbf{p\cdot x}}.
\end{gather*}
$f$ and $\tilde f$ can be set to create a ``wave packet'' of Gaussian form:
\[
    \tilde f(0,\mathbf{p})=\sigma^{-\frac{3}{4}}e^{-\frac{(\mathbf{p-k})^2}{2
    \sigma}}
\]
its anti-transform is
\[
    f(0,\mathbf{x})=\sigma^{\frac{3}{4}}e^{-\frac{x^2\sigma}{2}}e^{i\mathbf{k\cdot
    x}}.
\]
To compute their time evolution the Schrödinger picture is used:
\begin{gather*}
    i\pdv[]{}{t}\ket{t,f} = H\ket{t,f}\\
    \int_{}^{}\frac{\dd^3p}{2E_p}\pdv[]{\tilde f}{t}\ket{\mathbf{p}} =
    \int\frac{\dd^3p}{2E_p}\tilde f(t,\mathbf{p})\underbrace{H\ket{p}}_{
    =E_p\ket{\mathbf{p}}}\\
    i\pdv[]{\tilde f}{t} = E_p\tilde f(t,\mathbf{p} )
\end{gather*}
Imposing the initial condition stated before yields
\[
    \tilde f(t,\mathbf{p})=\sigma^{-\frac{3}{4}}e^{-\frac{(\mathbf{p-k})^2}{2
    \sigma}}e^{-iE_pt}
\]
The value of $E_p$ can be expanded around the value $p=k$, \textit{assuming}
$\sigma$ \textit{to be small},
\[
    E_p = E_k + \underbrace{\eval{\pdv[]{E_p}{p_i}}_{p=k}}_{:=v_k}(p_i-k_i)+...
\]
Plugging this in the exponent
\begin{gather*}
    f(t,\mathbf{x}) = \frac{1}{\sigma^{\frac{3}{4}}}\int \frac{\dd^3q}{2E_k}
    \exp\left[ -\frac{q^2}{2\sigma}-i\mathbf{v_k\cdot q}t - i\mathbf{q\cdot x}
    \right] \exp\left[ -iE_kt - i\mathbf{k\cdot x} \right]+...=\\=
    \frac{1}{\sigma^{\frac{3}{4}}} \int_{}^{} \frac{\dd^3q}{2E_k}
    \exp\left[ -\frac{1}{2\sigma}\left(\mathbf{q} + i\sigma(\mathbf{v_k}t+
    \mathbf{x})\right)^2 \right]+...
\end{gather*}
this is a simple Gaussian integral and yields\footnote{OCCHIO SEGNO DI X IN EXP}
\[
    f(t,\mathbf{x})=\sigma^{\frac{3}{4}}\frac{1}{2E_k}\exp\left[-\frac{\sigma}{2}
    \left(\mathbf{v_k}t-\mathbf{x}\right)^2 - iE_kt-i\mathbf{k\cdot x}\right]
\]
so this Gaussian is \textit{centred around its classical trajectory $x=v_kt$}.

For multi-particle states the form is similar:
\[
    \ket{f_1...f_n}=\int\prod_{i=1}^{n}\frac{\dd^3p_i}{2E_{p_i}}\tilde
    f_i(0,\mathbf{p_i}) \ket{\mathbf{p_1...p_n}}.
\]
The anti-transform at $t=0$ is the same as before for each particle:
\[
    f_i(0,\mathbf{x_i})=\sigma^{\frac{3}{4}}e^{-\frac{x_i^2\sigma}{2}}
    e^{i\mathbf{k_i\cdot x_i}}.
\]
And, unsurprisingly, also the time evolutions are the same
\begin{gather*}
    \tilde f_i(t,\mathbf{p_i})=\sigma^{-\frac{3}{4}}e^{-\frac{(\mathbf{p_i-k_i}
    )^2}{2\sigma}}e^{-iE_{p_i}t}\\
    f_i(t,\mathbf{x_i})=\sigma^{\frac{3}{4}}\frac{1}{2E_{k_i}}\exp\left[-\frac{
    \sigma}{2}\left(\mathbf{v_{k_i}}t-\mathbf{x_i}\right)^2 -
    iE_{k_i}t-i\mathbf{k_i\cdot x_i}\right]
\end{gather*}

\subsection{Interacting theory}
The fundamental assumption is that there exist a vacuum state ``close'' to the
free vacuum and that there exist a single particle state of momentum $\mathbf{p}$:
\[
    \ket{\mathbf{p}}\text{ with mass }m_p\sim m
\]
where $m$ is the bare mass (in AQFT it will be seen how $m$ and $m_p$ differ);
it is reasonable to assume this is an eigenstate of $H$ and of the momentum.

With this single particle state one can compute the quantities found before:
\begin{gather*}
    \ket{f} = \int_{}^{} \frac{\dd^3p}{2E_p}\tilde f(p)\ket{\mathbf{p}}\\
    \ket{t,f} = \frac{}{}\frac{\dd^3p}{2E_p}\tilde f(t,p)\ket{\mathbf{p}},
\end{gather*}
where $E_p = \sqrt{p^2-m_p^2}$.
To consider multiple particle states, one can assume that at $t=0$ they
are all in $x=0$, \textbf{but in the ``far'' past and future they do not interact,
therefore are free particle states}: this is the fundamental assumption
to consider scattering processes.

At $t=-\infty$ there is a state
\[
    \ket{\text{in}} = \int\prod_{i=1}^{n}\frac{\dd^3p_i}{2E_{p_i}}
    \tilde f_i(t,p_i)\ket{p_1...p_n}_\text{in}
\]
and similarly at $t=+\infty$
\[
    \ket{\text{out}} = \int\prod_{i=1}^{n}\frac{\dd^3p_i}{2E_{p_i}}
    \tilde f_i(t,p_i)\ket{p_1...p_n}_\text{out}.
\]

The free field
\[
    \phi(t,\mathbf{x}) = \int\dd^3p\left(g_p(t,x)a(p)+g^*_p(t,x)a^\dagger(p)
    \right) 
\]
where $g_p(t,x)$ can be found from the F.T. of the field.

It was shown in IQFT that
\begin{gather*}
    ia^\dagger (p)=\int_{}^{} \dd^3x g_p(t,x)\bar\pd_t\phi(t,x)\\
    -ia (p)=\int_{}^{} \dd^3x g^*_p(t,x)\bar\pd_t\phi(t,x)
\end{gather*}
where the derivatives act on both sides
\[
    g\bar\pd f = g\pd f - \pd g f.
\]
Their time derivatives are
\begin{gather*}
    i\pd_t a^\dagger (p) = \int_{}^{} \dd^3x\left(g_p(t,x)\ddot \phi - \ddot
    g_p(x,t)\phi\right) = \int_{}^{} \dd^3x \left(g_p(t,x)\ddot \phi - (
    \laplacian -m^2)g_p(t,x)\phi\right)
\end{gather*}
because $g_p$ satisfies the $K-G$ equation; this can be integrated by parts twice
to yield $(\Box +m^2) \phi=0$ in the free theory   
\[
    i\pd_ta^\dagger(p) = 0
\]
\section{LSZ reduction formula}
So in general these operators evolve in time, since for interacting fields the
$K-G$ equation is no longer homogeneous and the computed derivative does not
go to 0.
\begin{gather*}
    \lim\limits_{T\to-\infty} a^{(\dagger)}(t,p) = \sqrt{Z}a_\text{in}^{(\dagger)}
    (p)\\
    \lim\limits_{T\to+\infty} a^{(\dagger)}(t,p) =\sqrt{Z}a_\text{out}^{(\dagger)}
    (p)
\end{gather*}
[15 minuti]
